\documentclass{fdureport}
\usepackage{unicode-math}
\usepackage{longtable,booktabs,calc}
\providecommand{\tightlist}{\setlength{\itemsep}{0pt}\setlength{\parskip}{0pt}}
\usepackage{etoolbox}
\setlength{\emergencystretch}{3em}
\setcounter{secnumdepth}{0}
\AtBeginEnvironment{quote}{\begin{hintbox}}\AtEndEnvironment{quote}{\end{hintbox}}
\begin{document}
\def\fduimage{cover.jpg}\def\fducourse{功能学科综合实验}\def\fdutitle{实验题目（如：家兔高钾血症实验报告）}\def\fduexp{}\def\fduclass{年级专业}\def\fduname{姓名}\def\fduid{学号}\def\fdugroup{第
\_ 组}\def\fdudate{2025.\_\_.\_\_}\def\fduteacher{}\fducover
\fduheader{功能学科综合实验\quad 学号\quad 姓名}
\begin{center}{\heiti\Large\bfseries\color{ink} 实验题目（如：家兔高钾血症实验报告）}\\[2pt]{\small\kaishu\color{greyx} 年级专业\quad 姓名\quad 学号\quad 实验日期：2025.\_\_.\_\_}\end{center}
\vspace{-0.3em}{\color{gray!40}\hrule}\vspace{0.6em}
\begin{quote}
灰色引用块是\textbf{写作提示}，写完后请删除。结构参考课程常见要求，具体以任课老师当年的要求为准。
\end{quote}

\hypertarget{ux4e00ux5b9eux9a8cux76eeux7684}{%
\section{一、实验目的}\label{ux4e00ux5b9eux9a8cux76eeux7684}}

\begin{quote}
1--3
条，用动词开头（观察、理解、掌握）。写``这次实验要回答什么问题''，不是抄教材。
\end{quote}

\hypertarget{ux4e8cux5b9eux9a8cux539fux7406}{%
\section{二、实验原理}\label{ux4e8cux5b9eux9a8cux539fux7406}}

\begin{quote}
用自己的话写机制链：刺激 → 受体/通道 → 细胞/器官反应 →
可观测指标。必要时画一张示意图。
\end{quote}

\hypertarget{ux4e09ux5b9eux9a8cux5bf9ux8c61ux4e0eux65b9ux6cd5}{%
\section{三、实验对象与方法}\label{ux4e09ux5b9eux9a8cux5bf9ux8c61ux4e0eux65b9ux6cd5}}

\begin{quote}
动物（种类、体重、数量、麻醉方式与剂量）、药品（名称、浓度、给药途径）、仪器、关键操作步骤。\textbf{与实际做的不一致处要如实写}。
\end{quote}

\hypertarget{ux56dbux5b9eux9a8cux7ed3ux679c}{%
\section{四、实验结果}\label{ux56dbux5b9eux9a8cux7ed3ux679c}}

\begin{quote}
先文字概述，再放表和图。原始数据放表，趋势放图；每张图表有编号和标题，正文要引用。
\end{quote}

\textbf{表 1　（表题）}

\begin{longtable}[]{@{}lccc@{}}
\toprule\noalign{}
时间点 & 指标 1 & 指标 2 & 指标 3 \\
\midrule\noalign{}
\endhead
\bottomrule\noalign{}
\endlastfoot
对照 & & & \\
给药后 & & & \\
\end{longtable}

\hypertarget{ux4e94ux5206ux6790ux4e0eux8ba8ux8bba}{%
\section{五、分析与讨论}\label{ux4e94ux5206ux6790ux4e0eux8ba8ux8bba}}

\begin{quote}
建议结构：① 观察到了什么 → ② 与预期是否一致 → ③ 机制解释 → ④
与预期不一致的现象及可能原因（作为假设，不要当结论）→ ⑤
局限（样本量、操作误差）→ ⑥ 临床联系。
\end{quote}

\hypertarget{ux516dux7ed3ux8bba}{%
\section{六、结论}\label{ux516dux7ed3ux8bba}}

\begin{quote}
2--3 句，每句对应一个实验目的。
\end{quote}

\hypertarget{ux53c2ux8003ux6587ux732e}{%
\section{参考文献}\label{ux53c2ux8003ux6587ux732e}}

\begin{quote}
按课程要求的格式；教材：严钰锋，马淑兰
主编．功能学科实验教程［M］．2020．只写你真正读过的文献。
\end{quote}
\end{document}
